\begin{ExerciseBlock}
\ExerciseHeading{05}{Zug holt auf}{linear_kinematics_train_catches_up}
Ein Zug soll die Strecke $s=100 \,$km in der Zeit $T=120 \,$min zurücklegen.

\begin{enumerate}
	\item Welche Durchschnittsgeschwindigkeit $v_m$ muss der Zug fahren, dass er das Ziel in der vorgegebenen Zeit erreicht? Geben Sie ihre Ergebnis sowohl in km/h als auch in der SI-Basiseinheit an.

	\item Angenommen der Zug fährt die Hälfte der Zeit $T/2$ mit reduzierter Geschwindigkeit $v_1$. Mit welcher Geschwindigkeit $v_2$ muss er den zweiten Streckenabschnitt fahren um das Ziel noch in der vorgegebenen Zeit zu erreichen? Formulieren Sie ihr Ergebnis ohne Zahlenwerte einzusetzen. Werten Sie ihr Ergebnis für den Spezialfall $v_1 = v_m/2$ aus.

	\item Angenommen der Zug fährt die Hälfte der Strecke $s/2$ mit reduzierter Geschwindigkeit $v_1$. Mit welcher Geschwindigkeit $v_2$ muss er den zweiten Streckenabschnitt fahren um das Ziel noch in der vorgegebenen Zeit zu erreichen? Formulieren Sie ihr Ergebnis wieder ohne Zahlenwerte einzusetzen und werten Sie ihr Ergebnis für den Spezialfall $v_1 = v_m/2$ aus.

\end{enumerate}
\end{ExerciseBlock}
\UseExerciseSolution{linear_kinematics_train_catches_up}
